\section{Pre-Training}\label{sec:pretraining}

\begin{table}[t]
  \centering
  \caption{The \mossvl{} training curriculum: four pre-training stages
  followed by SFT and Realtime-SFT (\S\ref{sec:posttraining}). The data
  of each stage is described in the corresponding subsection. Token
  counts are the tokens fed to the language model, with vision tokens
  counted after the $2{\times}2$ token compression of the projection
  module (\S\ref{sec:architecture}); sequence lengths are in tokens.}
  \label{tab:stages}
  \small
\setlength{\tabcolsep}{6pt}
\renewcommand{\arraystretch}{1.15}
\begin{tabular}{@{}lrrrcc@{}}
\toprule
Stage & Tokens & Samples & Max seq. & Trainable & Peak LR \\
\midrule
\multicolumn{6}{@{}l}{\itshape Pre-training} \\
\quad 1\enspace Vision--language alignment & 150.3B & 219.4M & 8K   & Projection + cross-attn & $2{\times}10^{-4}$ \\
\quad 2\enspace Large-scale multimodal     & 203.0B & 139.3M & 64K  & Full model & $5{\times}10^{-5}$ \\
\quad 3\enspace High-quality multimodal    & 459.0B &  14.7M & 128K & Full model & $1{\times}10^{-5}$ \\
\quad 4\enspace Annealing \& long-context  & 450.1B &   9.3M & 256K & Full model & $1{\times}10^{-5}$ \\
\midrule
\multicolumn{6}{@{}l}{\itshape Post-training (\S\ref{sec:posttraining})} \\
\quad SFT          &  102.8B & 7.6M  & 128K & Full model & $1{\times}10^{-5}$ \\
\quad Realtime-SFT &   34.8B & 0.56M & 256K & Full model & $4{\times}10^{-5}$ \\
\bottomrule
\end{tabular}

\end{table}


\mossvl{} is pre-trained with a four-stage curriculum: vision--language
alignment, large-scale multimodal pre-training, high-quality multimodal
pre-training, and a final stage of annealing and long-context training.
Table~\ref{tab:stages} lists the token budget, sample count, maximum
sequence length, trainable modules, and peak learning rate of every
stage, including the two post-training stages of
\S\ref{sec:posttraining}. The table shows the shape of the
curriculum: the maximum sequence length grows from 8K to 256K tokens,
and the token budget shifts toward the later stages while sample counts
fall by orders of magnitude---many short samples early, far fewer but
much longer and denser ones late. Throughout, the data is
decontaminated against our evaluation suites.

A defining trait of this training run is the scale of our data
synthesis. Alongside data collected and reorganized from existing
corpora, we synthesize high-quality caption, OCR, grounding, and
temporal-grounding data at large scale throughout the curriculum. These
four types cover the perceptual fundamentals of a vision--language
model---describing scenes, reading embedded text, localizing objects,
and anchoring events in time---and this synthesized core underpins the
strong perception and temporal-understanding foundations of the
released models.

\subsection{Stage 1: Vision--Language Alignment}\label{sec:pt-align}

Stage 1 connects the two pre-trained components. Only the newly
introduced parameters---the projection module and the cross-attention
layers---are updated, while the vision encoder and the language model
remain frozen; the high peak learning rate in Table~\ref{tab:stages}
applies to these fresh modules alone. The data comprises two
categories, image captioning and OCR, and sequences stay short at 8K
tokens.

\subsection{Stage 2: Large-Scale Multimodal Pre-Training}\label{sec:pt-scale}

With the connectors aligned, Stage 2 unfreezes the full model and
supplies breadth. The mixture spans image and video captioning, OCR,
grounding, interleaved image--text documents, and text-only
pre-training corpora, together with multimodal understanding data over
single images, multi-image sets, videos, and plain text across diverse
domains, and reasoning data. The context window extends to 64K tokens,
which admits long interleaved documents and video.

\subsection{Stage 3: High-Quality Multimodal Pre-Training}\label{sec:pt-quality}

Stage 3 spends the largest token budget of the curriculum
(Table~\ref{tab:stages}) on its highest-quality data. The mixture keeps
the Stage-2 categories but rebalances them: captioning recedes,
multimodal understanding and reasoning data take a larger share, and
mathematics, knowledge-intensive data, and temporal grounding enter the
mixture. Sequences extend to 128K tokens.

\subsection{Stage 4: Annealing and Long-Context Training}\label{sec:pt-anneal}

The final stage combines long-context training with high-quality
annealing. One data strand consists of long-video captioning,
long-video QA, and long-video temporal grounding, together with
long-document and long-text data, mixed with a small share of the
regular categories, and stretches sequences to 256K tokens. The other is an annealing mixture that re-weights
toward mathematics, knowledge-intensive data, and instruction-tuning
and QA data, and includes identity data. The curriculum yields
\mossbase{}, the starting point for post-training
(\S\ref{sec:posttraining}).
